% Part: computability
% Chapter: tm-computations
% Section: church-turing-thesis

\documentclass[../../../include/open-logic-section]{subfiles}

\begin{document}

\olfileid{tur}{mac}{ctt}
\olsection{చర్చ్--ట్యూరింగ్ సిద్ధాంతప్రతిపాదన}

ప్రభావక ప్రక్రియ అనే భావనకు కచ్చితమైన ప్రత్యామ్నాయంగా ట్యూరింగ్ యంత్రాలు
ఉద్దేశించబడ్డాయి. ప్రభావక ప్రక్రియ భావననూ, ట్యూరింగ్ యంత్ర భావననూ గ్రహించిన
ఎవరైనా, ప్రభావక ప్రక్రియ ద్వారా చేయగల దేనినైనా ట్యూరింగ్ యంత్రంతో చేయవచ్చనే
అంతర్బోధను కలిగి ఉంటారని ట్యూరింగ్ భావించాడు. ప్రభావక ప్రక్రియ భావనకు
ప్రతిపాదించిన ఇతర కచ్చితమైన ప్రత్యామ్నాయాలన్నీ ట్యూరింగ్ యంత్ర భావనకు
విస్తారాత్మకంగా సమానమని తేలడం ఈ నిర్ధారణకు మద్దతిస్తుంది---అంటే అవన్నీ
కచ్చితంగా ఒకే ప్రమేయాల సమితిని గణించగలవు. ఈ నిర్ధారణను
\emph{చర్చ్--ట్యూరింగ్ సిద్ధాంతప్రతిపాదన} అంటారు.

\begin{defn}[చర్చ్--ట్యూరింగ్ సిద్ధాంతప్రతిపాదన]
ప్రభావక ప్రక్రియతో గణించగల ఏదైనా ట్యూరింగ్-గణనీయం అని
\emph{చర్చ్--ట్యూరింగ్ సిద్ధాంతప్రతిపాదన} చెబుతుంది.
\end{defn}

చర్చ్--ట్యూరింగ్ సిద్ధాంతప్రతిపాదనను రెండు విధాలుగా ఆశ్రయిస్తారు. మొదటి రకమైన
వాడకం బద్ధకానికి ఒక సాకు. ఏదో ఒక దాన్ని గణించే ప్రభావక ప్రక్రియకు, ఉదాహరణకు
``సూడోకోడ్''లో, మన వద్ద వివరణ ఉందనుకుందాం. అప్పుడు మనం నిజంగా ట్యూరింగ్
యంత్రాన్ని నిర్మించకపోయినా, అదే ప్రమేయాన్ని ఏదో ఒక ట్యూరింగ్ యంత్రం గణిస్తుందనే
నిర్ధారణను సమర్థించడానికి చర్చ్--ట్యూరింగ్ సిద్ధాంతప్రతిపాదనను ప్రయోగించవచ్చు.

చర్చ్--ట్యూరింగ్ సిద్ధాంతప్రతిపాదన యొక్క రెండవ వాడకం తాత్త్వికంగా మరింత
ఆసక్తికరమైనది. ట్యూరింగ్ యంత్రాలు గణించలేని ప్రమేయాలు ఉన్నాయని చూపవచ్చు.
దీనినుండి, చర్చ్--ట్యూరింగ్ సిద్ధాంతప్రతిపాదనను ఉపయోగించి, ఏ ప్రక్రియను
ఉపయోగించినా దాన్ని ప్రభావకంగా గణించలేమని నిర్ధారించవచ్చు. ఎందుకంటే అలాంటి
ప్రక్రియ ఉంటే, చర్చ్--ట్యూరింగ్ సిద్ధాంతప్రతిపాదన ప్రకారం దానికి ఒక ట్యూరింగ్
యంత్రం ఉండేదని అనుగమిస్తుంది. కాబట్టి దాన్ని గణించే ట్యూరింగ్ యంత్రం లేదని
నిరూపించగలిగితే, ప్రభావక ప్రక్రియ కూడా ఉండదు. ప్రత్యేకించి, ఆగే సమస్య అని
పిలిచేదాన్ని ట్యూరింగ్ యంత్రాలతో పరిష్కరించలేమని మాత్రమే కాక, అసలు ఏ
ప్రభావక పద్ధతిలోనూ పరిష్కరించలేమని చెప్పడానికి చర్చ్--ట్యూరింగ్
సిద్ధాంతప్రతిపాదనను ప్రయోగిస్తారు.

\end{document}
